\phantomsection
\chapter*{Simgeler ve Kısaltmalar}
\markboth{Simgeler ve Kısaltmalar}{Simgeler ve Kısaltmalar}
\addcontentsline{toc}{chapter}{Simgeler ve Kısaltmalar}

Bu bölümde kitap boyunca kullanılan temel simgeler ve kısaltmalar toplu olarak
verilmiştir. Büyük harfli $X$, $\bar X$ ve $S^2$ örneklem seçilmeden önce
rassal değişken veya tahmin ediciyi; küçük harfli $x$, $\bar x$ ve $s^2$ ise
gözlenen veriyi ve gerçekleşmiş sayısal değerleri gösterir. Evren
parametreleri Yunan harfleriyle, örneklem tahminleri Latin harfleri veya şapka
işaretiyle yazılır.

\small
\begin{longtable}{@{}>{\raggedright\arraybackslash}p{0.24\textwidth}>{\raggedright\arraybackslash}p{0.70\textwidth}@{}}
\toprule
\textbf{Simge} & \textbf{Anlamı} \\
\midrule
\endfirsthead
\toprule
\textbf{Simge} & \textbf{Anlamı} \\
\midrule
\endhead
\bottomrule
\endlastfoot
$N,n$ & Evren ve örneklem hacmi \\
\ifimoders
$k$ & Bağlama göre binom olay sayısı veya sistematik örnekleme aralığı \\
\else
$k$ & Bağlama göre binom olay sayısı, sistematik örnekleme aralığı veya ANOVA'daki grup sayısı \\
\fi
\ifimoders\else
$K$ & Regresyon modelindeki sabit terim dışındaki açıklayıcı sayısı \\
\fi
$\P(A)$ & $A$ olayının olasılığı \\
$X,Y$; $x,y$ & Rassal değişkenler ve gözlenen değerleri \\
$\E[X]$ & Beklenen değer \\
$\Var(X)$ & Varyans \\
$\mu,\sigma^2,p,\rho$ & Evren ortalaması, varyansı, oranı ve korelasyonu; $\rho$ küme içi korelasyonu da gösterebilir \\
$\bar X,S,S^2,\hat p$ & Örneklem ortalaması, standart sapması, varyansı ve oranı için tahmin ediciler \\
$\bar x,s,s^2,\hat p$ & Gözlenen örneklemde gerçekleşmiş tahminler \\
$\theta,\hat\theta$ & Hedef parametre ve bu parametrenin tahmin edicisi \\
$Q_1,Q_3,\IQR$ & Birinci ve üçüncü çeyrek ile çeyrekler arası açıklık \\
$\SE(\hat\theta)$ & Tahmin edicinin standart hatası \\
$\alpha$ & Anlamlılık düzeyi veya Tip I hata olasılığı \\
$\beta$ & Tip II hata olasılığı \\
\pval-değeri & Sıfır hipotezi altında gözlenen kadar veya daha aşırı sonuç elde etme olasılığı \\
$\Phi(z)$ & Standart normal dağılımın $z$ noktasındaki birikimli olasılığı \\
\ifimoders
$z_q,\ t_{q,\nu},\chi^2_{q,\nu}$ & İlgili dağılımların $q$ birikimli olasılık nicelikleri \\
\else
$z_q,\ t_{q,\nu},\chi^2_{q,\nu},F_{q;\nu_1,\nu_2}$ & İlgili dağılımların $q$ birikimli olasılık nicelikleri \\
\fi
$\nu$ & Serbestlik derecesi \\
\ifimoders
$Z,z,T,t,\chi^2$ & Standart normal, $t$ ve ki-kare test istatistikleri ile gerçekleşmiş değerleri \\
\else
$Z,z,T,t,\chi^2,F$ & Standart normal, $t$, ki-kare, ANOVA ve genel regresyon test istatistikleri ile gerçekleşmiş değerleri \\
\fi
$H_0,H_1$ & Sıfır ve alternatif hipotez \\
$\mu_0,\mu_1-\mu_2,\mu_D$ & Referans ortalama, iki evren ortalaması farkı ve eşleştirilmiş farkların evren ortalaması \\
$h$ & Güven aralığı planlamasında hedef hata payı \\
$D_i,\bar D,S_D$; $d_i,\bar d,s_D$ & Rassal fark nicelikleri ve gözlenen karşılıkları \\
$s_p,d,d_{\mathrm{tek}},d_z,g$ & Birleşik standart sapma ve kullanılan standartlaştırılmış fark ölçüleri \\
$O_{ij},E_{ij}$ & Çapraz tablodaki gözlenen ve beklenen hücre frekansları \\
$R_i,C_j$ & Çapraz tablonun satır ve sütun toplamları \\
$V$ & Cramér ilişki katsayısı \\
$\pi_i,w_i$ & $i$'inci birimin seçilme olasılığı ve örnekleme ağırlığı \\
$m,n_{\mathrm{etkin}},\DEFF$ & Ortalama küme büyüklüğü, etkin örneklem hacmi ve tasarım etkisi \\
\ifimoders\else
$B,\hat\theta_b^*$ & Yeniden örnekleme tekrar sayısı ve $b$'inci bootstrap tahmini \\
$\bar\theta^*,T_{\mathrm{göz}},T_b^*$ & Bootstrap tahminlerinin ortalaması, gözlenen test istatistiği ve $b$'inci rastgeleleştirme istatistiği \\
$\Ind{A}$ & $A$ koşulu doğruysa 1, değilse 0 alan gösterge fonksiyonu \\
$\SE_{\mathrm{boot}},p_{\mathrm{rand}}$ & Bootstrap standart hatası ve rastgeleleştirme testi $p$ değeri \\
\fi

$r_{ij},r_{ij}^{\ast}$ & Pearson hücre artığı ve düzeltilmiş standartlaştırılmış artık \\
$r$ & Örneklem Pearson korelasyonu \\
\ifimoders
$\beta_0,\beta_1,b_0,b_1$ & Basit regresyonda evren sabiti ve eğimi ile örneklem tahminleri \\
\else
$\beta_j,b_j$ & $j$'inci evren regresyon katsayısı ve örneklem tahmini \\
\fi
\ifimoders
$\varepsilon_i,e_i,\hat Y,\hat y_i$ & Model hatası, gözlenen artık, uydurulan doğru ve gerçekleşmiş uydurulan değer \\
\else
$\varepsilon_i,e_i,\hat Y_i,\hat y_i$ & Model hatası, gözlenen artık, uydurulan değer kuralı ve gerçekleşmiş uydurulan değer \\
\fi
$R^2$ & Regresyon modelinin açıklanan varyans oranı \\
\ifimoders\else
$R^2_{\mathrm{düz}}$ & Açıklayıcı sayısını dikkate alan düzeltilmiş açıklanan varyans oranı \\
$\SSE,\SST$ & Hata kareleri toplamı ve toplam kareler toplamı \\
$R_j^2,\VIF_j$ & $X_j$ için yardımcı regresyonun açıklanan varyans oranı ve varyans büyütme faktörü \\
$SS_{\mathrm{gruplar}},SS_{\mathrm{hata}}$ & ANOVA'da gruplar arası ve grup içi kareler toplamları \\
$SS_{\mathrm{toplam}}$ & ANOVA'da toplam kareler toplamı \\
$MS_{\mathrm{gruplar}},MS_{\mathrm{hata}}$ & ANOVA'da gruplar arası ve grup içi ortalama kareler \\
$\eta^2,\omega^2$ & ANOVA için eta-kare ve omega-kare etki büyüklükleri \\
\fi

\end{longtable}

\vspace{0.5em}

\begin{longtable}{@{}>{\raggedright\arraybackslash}p{0.24\textwidth}>{\raggedright\arraybackslash}p{0.70\textwidth}@{}}
\toprule
\textbf{Kısaltma} & \textbf{Açılımı} \\
\midrule
\endfirsthead
\toprule
\textbf{Kısaltma} & \textbf{Açılımı} \\
\midrule
\endhead
\bottomrule
\endlastfoot
\ifimoders\else
ANOVA & Analysis of variance; varyans analizi \\
\fi
\ifimoders\else
BCa & Yanlılık düzeltilmiş ve hızlandırılmış bootstrap aralığı \\
\fi
CI & Confidence interval; İngilizce yazılım çıktılarında güven aralığı \\
DEFF & Tasarım etkisi \\
df & Degrees of freedom; yazılım kodunda serbestlik derecesi \\
GA & Güven aralığı \\
GAISE & İstatistik Eğitimi için Değerlendirme ve Öğretim İlkeleri \\
\ifimoders\else
HSD & Honestly significant difference; Tukey çoklu karşılaştırma yöntemi \\
\fi
IQR & Çeyrekler arası açıklık \\
MLT & Merkezi Limit Teoremi \\
MSE & Ortalama karesel hata \\
\ifimoders\else
OLS & En küçük kareler yöntemi \\
\fi
SE & Standart hata \\
sd & Serbestlik derecesi \\
SPSS & Statistical Package for the Social Sciences \\
\ifimoders\else
VIF & Varyans büyütme faktörü \\
\fi
\end{longtable}
\begin{prismbox}[PrismLight]{Yazım ve raporlama standardı}
Metinde \pval-değeri biçimi kullanılır. Güven aralığı Türkçe metinde
\textbf{GA}, İngilizce yazılım çıktısı açıklanırken gerekirse \textbf{CI}
olarak gösterilir. Serbestlik derecesi formüllerde $\nu$, Türkçe raporlama
cümlelerinde ``sd'' ve yazılım kodunda ilgili işlev gerektiriyorsa
\texttt{df} ile belirtilir. Kuramsal rassal nicelikler büyük, gözlenen
değerler küçük harfle yazılır.
\end{prismbox}
\normalsize